\documentclass[11pt]{article}

\usepackage[a4paper,margin=1in]{geometry}
\usepackage{amsmath,amssymb,amsthm,bm}
\usepackage{booktabs,longtable,tabularx,array}
\usepackage{enumitem}
\usepackage{xcolor}
\usepackage{hyperref}
\usepackage{fancyhdr}
\usepackage{microtype}
\usepackage{listings}
\usepackage{caption}

\hypersetup{colorlinks=true,linkcolor=blue,urlcolor=blue,citecolor=blue}
\setlength{\headheight}{14pt}
\pagestyle{fancy}
\fancyhf{}
\lhead{CoPE vs.\ FSR-PC: Experiment Design}
\rhead{Internal Handoff v2}
\cfoot{\thepage}

\definecolor{codegray}{rgb}{0.95,0.95,0.95}
\definecolor{warnbg}{rgb}{1.0,0.96,0.90}
\lstset{backgroundcolor=\color{codegray},basicstyle=\ttfamily\footnotesize,
  breaklines=true,frame=single,columns=fullflexible}

\newcommand{\cope}{\textsc{CoPE}}
\newcommand{\fsrpc}{\textsc{FSR-PC}}
\newcommand{\noadapt}{\textsc{no-adaptation}}

\title{\bfseries Persistent Local Constraint-State Patching versus
Full-State Regeneration\\[4pt]
\large Experiment design, fairness protocol, and CPU pilot results}
\author{Internal technical report --- CPU phase}
\date{2026--08--03}

\begin{document}
\maketitle

\begin{abstract}
We compare two designs for adapting a robot's task state when a user interrupts
a long-horizon task: \cope{} (Constraint-state Patch Editing), which edits a
persistent constraint state with minimal typed patches, and \fsrpc{}, an
\emph{internally defined} strong baseline that regenerates the complete state at
every interruption. The comparison is run on a purpose-built multi-object
basket-sorting benchmark whose nominal reliability is gated at
$\geq 8/10$ before any method is compared on it, so that adaptation failure is
not confounded with execution failure. This report documents the method, the
baseline's provenance, the benchmark, the fairness controls, the metrics, and
the CPU pilot. Both arms enter one shared downstream stack --- the frozen
online-repair engine --- so that only their state-update semantics differ; every
interruption is shown to drive continuation capture, runtime operator synthesis,
sequential rollout verification, hard rejection, graph splicing, restore
validation and stage resumption. On the CPU pilot the two arms are \emph{tied on
every behavioural outcome}; the measurable differences are in state
representation, edit locality, and which audit questions the decision record can
answer. No robotics claim is made: the development machine is CPU-only and no
GPU software was installed, imported, or executed.
\end{abstract}

\tableofcontents

\section{Scope, and what is deliberately not claimed}

\begin{enumerate}[leftmargin=*]
\item The earlier ReKep online-repair work is retained as an \textbf{integration
case study only}. Its base task has a low nominal success floor and confounds
execution failure with recovery failure, which makes it unsuitable as the
primary benchmark for a question about task-state adaptation.
\item The existing online repair core is \textbf{reused, not rewritten}. The
\cope{} layer sits above it and an adapter maps between the two
representations; nothing in the frozen engine was replaced or duplicated.
\item \fsrpc{} is \textbf{not an established published algorithm}
(Section~\ref{sec:fsrpc}). It must never be cited as prior art.
\item All numbers in this report come from the \textbf{frozen 2D synthetic
environment} on CPU. Rollout verification is therefore meaningful --- the
verifier simulates candidates under the same kinematics and control limits the
executor uses --- but this is still \textbf{not a robotics result}.
\item \cope{} is \textbf{not claimed} to beat \fsrpc{} on final task success. On
the pilot the two arms tie at $20/20$.
\end{enumerate}

\section{The \cope{} semantic layer}

Terminology follows the collaborator's memo (Jingsu Li, \emph{CoPE / RCSP
Comprehensive Technical Memo}, v0.3, 2026--04--26): \cope{} is
\emph{Constraint-state Patch Editing} and RCSP is \emph{Reactive
Constraint-State Patching}. No expansion used here is invented.

\subsection{Slot schema}
A constraint is a software-engineering object, not only a mathematical function:
\[
c_i = (\mathrm{id},\; \omega,\; \mathrm{mode},\; \mathrm{priority},\;
\mathrm{source},\; \mathrm{validity},\; \mathrm{restore},\; \mathrm{lineage},\;
\mathrm{history}),
\]
where the grounding $\omega$ is one field among nine.

\subsection{Lifecycle: existence separated from participation}
\begin{center}
\begin{tabular}{@{}llll@{}}
\toprule
mode & exists & participates in planning & restorable\\
\midrule
\texttt{active} & yes & yes & ---\\
\texttt{suspended} & yes & no & yes\\
\texttt{overridden} & yes & no & yes\\
\texttt{expired} & yes (auditable) & no & \textbf{never}\\
\bottomrule
\end{tabular}
\end{center}
\texttt{expired} is a soft delete: the record survives so that ``which goal was
cancelled, and why?'' stays answerable. Slots are never removed.

\subsection{Typed patch operators}
The first-paper subset is
$\{\textsf{Insert},\textsf{Suspend},\textsf{Override},\textsf{Inherit},
\textsf{Expire},\textsf{Revalidate},\textsf{Restore}\}$. Two rules are pinned
down by tests because violating them would still produce a passing benchmark
number:
\begin{enumerate}[leftmargin=*]
\item \textsf{Revalidate} is a \emph{pure check}: it returns
\textsf{OK}/\textsf{DEGRADED}/\textsf{FAIL} and never changes a mode.
\item \textsf{Restore} is a \emph{commitment}: it is refused unless revalidation
returns \textsf{OK}, refused outright for expired slots, and a \textsf{Restore}
not preceded by a \textsf{Revalidate} is a validation violation rejected before
any mutation.
\end{enumerate}

\subsection{Invariants}
Checked before and after every patch: identity stability, lifecycle legality
(\texttt{expired} terminal), no deletion, symmetric and acyclic override edges,
per-slot history monotonicity, and guarded restoration. A patch failing
pre-validation is not applied at all; the state is left untouched.

\subsection{Two layers}
\begin{center}
\begin{tabular}{@{}p{0.44\textwidth}p{0.48\textwidth}@{}}
\toprule
\cope{} semantic layer (this work) & repair execution layer (frozen, reused)\\
\midrule
persistent state $S_t$, relation graph $G_t$, history $H_t$, typed patch
operators &
continuation capture $\to$ operator synthesis $\to$ legality $\to$ sequential
rollout verification $\to$ splice $\to$ restore $\to$ resume\\
\bottomrule
\end{tabular}
\end{center}
The adapter maps \texttt{ConstraintSlot(mode=SUSPENDED)} to a stage suspension,
a \texttt{Patch} to a repair goal, and a slot's restore predicate plus lineage
to a continuation. The repair engine remains the sole owner of execution,
verification, and splicing.

\section{\fsrpc{}: provenance and specification}
\label{sec:fsrpc}

\fcolorbox{black}{warnbg}{\parbox{0.97\textwidth}{%
\textbf{\fsrpc{} is an internally defined baseline.} The acronym does not appear
in the collaborator's memo or in any other supplied material; a repository-wide
search returned zero matches before this work began. The expansion used
throughout, \emph{Full-State Regeneration with Progress/Context}, comes only
from the project brief that commissioned the baseline. Describe it as ``our
internally defined strong baseline''. Do not attach a citation. Do not imply
prior art.}}

\medskip
At each interruption,
\[
\text{\cope{}: } S_{t+1} = \mathrm{apply}(P_t, S_t),\quad P_t = [\mathrm{op}_1
\ldots \mathrm{op}_k],\ k \text{ small};\qquad
\text{\fsrpc{}: } S_{t+1} = \mathrm{regenerate}(\mathrm{context}_t).
\]

\paragraph{Why it is strong, not a straw man.}
\fsrpc{} receives byte-identical input to \cope{} at every interruption --- a
claim asserted by a test over all four conditions and three seeds, not by prose.
The shared packet carries the full world and robot state, completed goals,
pending and cancelled goals, the natural-language user request, all safety
constraints, the complete task and event histories, perception, and the same
compute budget and horizon. The implementation honours completed goals rather
than re-planning them, carries prior lifecycle modes forward, collapses its own
re-minted slots onto the canonical task vocabulary so repeated regenerations do
not multiply records, and has access to the same naming helper \cope{} uses.
It is explicitly \emph{not} made forgetful.

\paragraph{The single intended difference.}
\begin{center}
\begin{tabular}{@{}lll@{}}
\toprule
 & \cope{} & \fsrpc{}\\
\midrule
adaptation output & typed patch over existing slots & complete new state\\
slot identity & carried & re-minted\\
per-slot edit history & continued & not inherently continued\\
lineage / override edges & explicit & not represented\\
edit footprint & touched subset & whole state\\
reason attached to a change & per operation & none\\
\bottomrule
\end{tabular}
\end{center}

\section{Benchmark: multi-object basket sorting}

Three objects (milk, yogurt, butter) and three baskets. Nominal assignment:
milk$\to$A, yogurt$\to$A, butter$\to$B. \texttt{basket\_C} exists as a spare
destination so that a \emph{repeated} redirection can end somewhere other than
the nominal target; without it, condition I3 would be a no-op for final
placement and the control arm would pass it for free. Two hard safety
constraints are always present. Execution follows a fixed scripted order carried
in the shared slot payload.

\paragraph{Reliability gate (Stage A): a precondition, not a result.}
Nominal task success must reach $\geq 8/10$ over ten seeds ($\geq 9/10$
preferred), measured with the \noadapt{} arm on the uninterrupted task so that
the number describes the task and the executor rather than any method.
\textbf{If the gate fails, the task is changed or simplified; the methods are
not tuned.} The CPU result is $10/10$, and the gate is verified non-vacuous: a
degraded executor ($p=0.5$) fails it.

\paragraph{Interruption conditions.}
\begin{description}[leftmargin=1.4em,style=nextline]
\item[I1 --- cancellation + redirection] after milk: cancel yogurt, redirect
butter B$\to$A.
\item[I2 --- temporary unavailability] before any goal: basket B disappears, so
the butter goal must be set aside and later brought back.
\item[I3 --- repeated interruption] the same goal is redirected twice
(B$\to$A, then A$\to$C), plus a new safety constraint. Both updates land before
butter is attempted, so the condition measures state survival rather than
reaction latency.
\item[I4 --- changed-world restoration] as I2, but basket B is \emph{moved}
while unavailable, so restoring the suspended goal requires revalidating its
grounding rather than restoring blindly.
\end{description}

\section{Integration with the frozen repair engine}

An internal audit of the v1 package found that the CoPE semantic layer was
implemented but the online-repair execution core was neither delivered nor
exercised: v1 tested \emph{event $\to$ patch/regeneration $\to$ active-slot
compilation $\to$ mock executor}. Every finding was verified and correct. This
version integrates the engine.

\paragraph{What runs now.}
\[
\begin{aligned}
C_t &= (S_t, G_t, H_t), &
P_t &= \epsilon(e_t, C_t, x_t), &
C_t^{+} &= \mathrm{Apply}(C_t, P_t),\\
\kappa_t &= \mathrm{Capture}(\mathrm{TP}, \text{stage}, x_t), &
\Omega_t &= \mathrm{Generate}(C_t^{+}, x_t, \kappa_t), &
\Delta_t^{*} &= \arg\min_{\Delta} J(\Delta),
\end{aligned}
\]
subject to \textsc{Legal}, \textsc{RolloutFeasible} and \textsc{RestoreValid},
followed by $\mathrm{TP}^{+} = \mathrm{Splice}(\mathrm{TP}, \Delta_t^{*},
\kappa_t)$ and Execute $\to$ Revalidate $\to$ Restore $\to$ Resume.

Everything from $\mathrm{Capture}$ onwards is the frozen engine, imported and
called: \texttt{TaskProgram}, \texttt{Continuation}, \texttt{RepairPlanner},
\texttt{SynthesisGenerator}, the legality and feasibility filters,
\texttt{RolloutVerifier}, \texttt{Scorer}, the splice, the restore gate, the
controllers and guards, and \texttt{Synthetic2DEnv}. The repair core is
byte-identical to the pre-pivot commit.

\paragraph{The bridge.}
The genuinely new component is a \emph{constraint-state to repair-goal
compiler}. It reads the state \emph{diff} --- never the policy --- and maps each
lifecycle transition onto the engine's own goal vocabulary: a re-grounded slot
demands \texttt{aligned\_to\_current\_target}; a suspended, cancelled or
restored slot demands \texttt{restore\_valid}; an inserted hard safety slot
demands \texttt{safe\_clearance}. Every requirement carries the slot and
transition that induced it, so the trace answers \emph{why} the robot had to act,
not only \emph{what} it did.

\paragraph{Verification is no longer vacuous.}
Against a symbolic mock, a rollout verifier has nothing to predict. Executing
legs through the frozen environment gives it the same kinematics and control
limits the executor uses, so a candidate that would collide, time out or fail to
restore is genuinely rejected before execution. Two modelling errors were caught
this way and corrected: a suspended goal was wrongly required to clear an
obstruction, and an unavailable basket was wrongly modelled as a permanent
obstacle. In both cases the verifier rejected every candidate --- correct
behaviour applied to an incorrectly stated problem.

\paragraph{Evidence.} Over the 60-episode pilot: 105 repair invocations, 165
candidate programs synthesised at runtime, 165 sequentially rollout-verified,
105 graph splices, 105 restore validations, 105 stage resumptions, 0
safe-fallbacks, and 0 primary-arm episodes with an incomplete pipeline. A gate
script prints the ten-marker sequence per interruption and exits non-zero if any
is incomplete.

\paragraph{Both arms share the stack.} A single
\texttt{enter\_repair\_pipeline} entry point is called by both policies with the
identical signature. Measured across conditions and seeds, the arms receive
identical adaptation inputs \emph{and} compile to identical repair goals, and
every paired bootstrap interval on downstream work (repairs, candidates,
rollouts, splices, restores, resumes) is exactly $[0,0]$.

\section{Fairness protocol}

Every control is implemented in one place --- the shared episode runner --- so
it cannot drift between arms, and each is asserted by a test.

\begin{description}[leftmargin=1.4em,style=nextline]
\item[F1 same task] the pairing key $(\text{task},\text{condition},\text{seed})$
deliberately excludes the method; the initial state is built before any policy
object exists.
\item[F2 same information] each arm receives one adaptation input built from
method-independent facts; the serialized packet is fingerprinted, and the
fingerprints are byte-identical at every event for all conditions and seeds.
\item[F3 same executor] both arms emit the same compiled request structure from
the same state store; the executor object and its random stream are shared, and
the stream is keyed on the \emph{physical action} so the arms get common random
numbers. The compiled request contains no method marker.
\item[F4 same budget, horizon, grader, metric code] the grader depends only on
the condition and never inspects what a policy did.
\item[F5 no cross-contamination] episodes are independent; store, trace, and
executor are fresh per episode.
\item[F6 one shared repair engine] a single \texttt{SharedRepairEngine} instance
serves whichever policy is running, so the arms cannot diverge downstream even
by accident.
\end{description}

\paragraph{One legitimate divergence, stated explicitly.}
The \noadapt{} control can see a different \emph{second} packet under I2/I4,
because it really did something different first: it burned an action on an
unavailable target. Its \emph{first} packet is identical to the others, which is
the information test. For the comparison the paper makes --- \cope{} vs
\fsrpc{} --- packets are identical at every event.

\paragraph{Statistics.}
Exact McNemar on discordant pairs with Wilson 95\% intervals for binary
outcomes; paired bootstrap (10\,000 resamples, seed 0) for continuous outcomes;
Holm--Bonferroni across the binary family. The statistics module is reused
unchanged from the frozen CPU work.

\section{Metrics}

Five families, all computed by one function from the shared store, trace, and
executor log: final behaviour; progress and continuity; state representation;
editing versus regeneration; and audit. \texttt{None} always means \emph{not
applicable to this episode}, never \emph{failed}.

\paragraph{Audit, and why it is not circular.}
The record must answer: Q1 which goal was cancelled; Q2 which completed goal was
preserved; Q3 why was a constraint suspended; Q4 what condition allowed
restoration; Q5 which state replaced the old target. A trace made only of
\cope{} operations would answer all five by construction and score \fsrpc{} at
zero for circular reasons. Every query therefore has \emph{two} evidence paths:
typed patch operations, \emph{and} diffing consecutive regenerated snapshots.
Q1 and Q2 are recoverable from a state snapshot, and \fsrpc{} answers them.
Q3--Q5 ask for a reason, a restoration condition, and a replacement identity ---
none of which a complete snapshot carries. That asymmetry is the finding.
Queries a condition does not pose score \texttt{None}, never \texttt{False}.

\section{CPU pilot results}

Stage A: $10/10 \to$ PASS. Stage B: 60 episodes (5 paired seeds $\times$ 4
conditions $\times$ 3 arms).

\begin{table}[h]
\centering
\caption{Binary outcomes, exact McNemar on 20 paired episodes per arm.}
\begin{tabular}{@{}lccc@{}}
\toprule
metric & \cope{} & \fsrpc{} & $p$\\
\midrule
revised task success & 20/20 & 20/20 & 1.000\\
completed progress preserved & 20/20 & 20/20 & 1.000\\
remaining goal correct & 20/20 & 20/20 & 1.000\\
\textbf{identity preserved} & \textbf{20/20} & \textbf{0/20} & $1.9\times10^{-6}$\\
lifecycle transitions legal & 20/20 & 20/20 & 1.000\\
lineage correct & 20/20 & 20/20 & 1.000\\
history monotonic & 20/20 & 20/20 & 1.000\\
\bottomrule
\end{tabular}
\end{table}

\begin{table}[h]
\centering
\caption{Continuous outcomes, paired bootstrap 95\% CI on
$\text{\cope{}}-\text{\fsrpc{}}$.}
\begin{tabular}{@{}lcccc@{}}
\toprule
metric & \cope{} & \fsrpc{} & 95\% CI & excludes 0\\
\midrule
normalized edit distance & 0.457 & 1.000 & $[-0.611,-0.472]$ & yes\\
slots touched per event & 2.75 & 5.13 & $[-2.850,-1.875]$ & yes\\
\textbf{audit coverage} & \textbf{1.000} & \textbf{0.458} & $[+0.483,+0.600]$ & yes\\
completion steps & 44.5 & 44.5 & $[0,0]$ & no\\
invalid actions & 0.000 & 0.000 & $[0,0]$ & no\\
repairs / candidates / splices / resumes & \multicolumn{2}{c}{identical} & $[0,0]$ & no\\
\bottomrule
\end{tabular}
\end{table}

\begin{table}[h]
\centering
\caption{Per-condition final success on the physical stack. Every interruption
condition now punishes not adapting.}
\begin{tabular}{@{}lccc@{}}
\toprule
condition & \cope{} & \fsrpc{} & \noadapt{}\\
\midrule
I1 & 5/5 & 5/5 & \textbf{0/5}\\
I2 & 5/5 & 5/5 & \textbf{0/5}\\
I3 & 5/5 & 5/5 & \textbf{0/5}\\
I4 & 5/5 & 5/5 & \textbf{0/5}\\
\bottomrule
\end{tabular}
\end{table}

\paragraph{Stage C ablation: mechanism, not rhetoric.}
\begin{center}
\begin{tabular}{@{}lcccc@{}}
\toprule
arm & success & identity preserved & mean edit distance & mean audit coverage\\
\midrule
\cope{} & 20/20 & 20/20 & 0.507 & 1.000\\
\fsrpc{} & 20/20 & 0/20 & 1.000 & 0.458\\
\fsrpc{}(preserve ids) & 20/20 & \textbf{20/20} & 1.000 & 0.458\\
\bottomrule
\end{tabular}
\end{center}
Re-using identifiers recovers the identity metric and changes nothing else. The
locality and audit differences therefore come from local typed editing, not from
identifier hygiene.

\section{Honest interpretation}

\paragraph{Supported.}
The arms tie on every behavioural outcome. The measurable differences are
structural: identity survival, edit locality, and which audit questions the
record can answer. The audit gap is the substantive one, and the two-evidence-path
scoring is what makes it non-circular.

\paragraph{Not supported.}
A normalized edit distance of $1.000$ for \fsrpc{} is true \emph{by definition}
of full regeneration; it describes the two designs and must not be presented as
an experimental finding. No physical, robotics, or wall-clock-cost claim
follows from a CPU mock. Neither arm uses a language model, so no token or
prompt-cost claim is made. With 20 paired episodes per arm, the behavioural ties
mean \emph{no evidence of a difference}, not \emph{evidence of no difference}.

\paragraph{What would change the conclusion.}
The mock executor never punishes a slightly-wrong-but-recoverable state. On the
GPU host, the arms may separate under longer horizons where regenerated state
drifts across many interruptions; under more than two interruptions, where
re-minted identity obscures what has already been attempted; or where the
executor cannot cheaply retry, so a duplicated or stale goal becomes
unrecoverable. If none of these separate the arms, the honest paper claim is
about representation, auditability, and edit locality --- not task success ---
and the results section should say exactly that.

\section{Status and next steps}

Local validation complete: 229 tests pass (94 pre-existing, unchanged; 135 new),
Stage A/B/C pilots run, five side-by-side dry-run traces produced, this report
compiles, and checksums are recorded. Nothing GPU-related has been executed;
every GPU command in the runbook is marked \textsc{remote gpu server only} and
is untested by construction.

The single integration point for the GPU host is the executor class. Both
policies already emit the same compiled request; no policy code changes.

\section*{Authorship and attribution}

The abstract idea of online constraint-program repair, and the CoPE/RCSP
constraint-state formulation with its slot schema, lifecycle modes, and typed
patch operators, originate with the collaborator (Jingsu Li, memo v0.3,
2026--04--26). The work documented here is the experimental realisation: the
adapter onto the existing repair engine, the internally defined \fsrpc{}
baseline, the benchmark, the fairness protocol, the metrics, and the CPU pilot.
Author order and final framing are not settled here and should not be inferred
from this document.

\end{document}
